\documentclass{situationist-preprint}

\title{The Arity of Givenness}
\msid{PHE-02-a}
\msfamily{PHE-02}
\msdiscipline{Phenomenology}
\mstype{Principal manuscript}
\mskeywords{givenness; intentionality; as-structure; myth of the given; Marion; Henry; Heidegger}

\begin{document}
\maketitle

\begin{abstract}
What is it for something to be given? We argue that givenness is an irreducibly three-place relation: something is given to someone as something. Each of the three places has been the target of an influential elimination. Sense-datum theories and Marion's pure givenness drop the \enquote{as}, treating the given as prior to every determination; Heidegger's late \enquote{it gives} and some readings of Marion drop the dative; Henry's self-affection drops the something, making givenness give only itself. We state a criterion of elimination, on which a place is eliminable if the truth of givenness claims is fixed by the other two places, and show by three families of minimal pairs that no place meets it: the same thing is given to the same subject as different things, the same thing is given as the same to one subject and not to another, and different things are given to one subject as the same. We then argue that each elimination, examined closely, either reintroduces the missing place or changes the subject from givenness to something else: causation, being, or life.
\end{abstract}

\section{The question}

Phenomenology is the study of things as they are given, and the word \enquote{given} carries most of its methodological weight. Husserl's principle of all principles states that whatever presents itself in intuition in its originary form is to be accepted as it is given, within the limits in which it is given \citep[§24]{Husserl1913}. Yet what it is for something to be given is rarely stated in a form that allows its logical structure to be inspected. The question we ask is accordingly a question of arity: how many places does the relation have, and can any of them be eliminated?

Three places appear in the ordinary statement. Something (the cube, the melody, the other's anger) is given; it is given to someone (a perceiver, a listener, an empathizer); and it is given as something (as a cube, as a melody heard from the middle, as anger rather than fatigue). Husserl's analysis of intentionality already contains all three: the object intended, the act and its ego, and the noematic sense through which the object is intended \citep[§§88--91]{Husserl1913}. Heidegger's \enquote{as-structure} of understanding makes the third place explicit: what is encountered is always encountered as a table, as a door, as something to be used \citep[§32]{Heidegger1927}.

Each place, however, has been contested. The empiricist tradition and its critics treated the given as what is present before interpretation, and \citet{Sellars1956} argued that nothing could be both given in that sense and capable of justifying belief. \citet{Marion1989,Marion1997} argued that the reduction, pushed far enough, reaches a givenness that is prior to every horizon and every concept, and in his formula, the more reduction, the more givenness. \citet{Heidegger1969} sought to think the \enquote{it gives} in \enquote{there is being} without a giver and, on a natural reading, without a recipient in the sense of a subject. \citet{Henry1963} argued that the original form of manifestation is self-affection, in which life is given to itself and to nothing else. Each proposal deletes a place. We argue that none succeeds.

\section{Framework}

\begin{definition}\label{d:giv}
A \emph{givenness claim} has the form $G(x, s, F)$: the item $x$ is given to the subject $s$ as $F$, where $F$ is a sense, that is, a way the item is presented (as a cube, as seen from the left, as threatening). Givenness claims are evaluated at a time, which is held fixed in what follows.
\end{definition}

\begin{definition}\label{d:elim}
A place of $G$ is \emph{eliminable} if there is a relation $R$ on the remaining two places such that, for all values of the three places, $G$ holds exactly when $R$ holds of the remaining two.
\end{definition}

Definition~\ref{d:elim} is the weakest notion of elimination that the eliminative proposals need. If a place cannot be eliminated in this sense, it cannot be eliminated by any more demanding analysis either.

\begin{principle}\label{p:aspect}
One and the same item can be given to one and the same subject under different senses.
\end{principle}

\begin{principle}\label{p:subject}
One and the same item can be given under one and the same sense to one subject and not to another.
\end{principle}

\begin{principle}\label{p:item}
Different items can be given to one and the same subject under one and the same sense.
\end{principle}

The three principles are descriptions of familiar experiences. For Principle~\ref{p:aspect}: the duck-rabbit figure is given to one viewer first as a duck and then as a rabbit, without change in the figure or in the viewer \citep[Part II, §xi]{Wittgenstein1953}; the man at the door is given to me as a stranger and, a moment later, as my brother returned after many years. For Principle~\ref{p:subject}: the cube on the table is given as a cube to the person facing it and not to the person in the next room; the patient's pain is given as pain to the patient primordially and, if PHE-12 is right, to the physician non-primordially, which is a difference in the subject's mode of access. For Principle~\ref{p:item}: two qualitatively identical cups, swapped while the drinker looks away, are given to her as the same cup, her cup.

\section{The theorem}

\begin{lemma}\label{l:as}
The sense place is not eliminable.
\end{lemma}

\begin{proof}
Suppose there is $R$ on items and subjects with $G(x,s,F)$ iff $R(x,s)$ for all $x$, $s$, $F$. By Principle~\ref{p:aspect}, there are $x$, $s$, $F$ and $F'$ with $G(x,s,F)$ and not $G(x,s,F')$: at the moment the figure is given as a duck, it is not given as a rabbit. Then $R(x,s)$ holds and fails, a contradiction.
\end{proof}

\begin{lemma}\label{l:dative}
The subject place is not eliminable.
\end{lemma}

\begin{proof}
By Principle~\ref{p:subject}, there are $x$, $F$, $s$ and $s'$ with $G(x,s,F)$ and not $G(x,s',F)$. A relation $R(x,F)$ cannot agree with both.
\end{proof}

\begin{lemma}\label{l:item}
The item place is not eliminable.
\end{lemma}

\begin{proof}
By Principle~\ref{p:item}, take the swapped cups: after the swap, the cup $x'$ before the drinker $s$ is given to her as her cup ($F$), and the original cup $x$, now in the cupboard, is not given to her at all, so not as $F$. A relation $R(s,F)$ cannot agree with both $G(x',s,F)$ and the failure of $G(x,s,F)$.
\end{proof}

\begin{theorem}\label{t:main}
Givenness is an irreducibly three-place relation: something is given to someone as something, and no place can be eliminated.
\end{theorem}

\begin{proof}
By Lemmas~\ref{l:as}, \ref{l:dative} and~\ref{l:item}.
\end{proof}

The proofs are short because the work is done by the principles, and the principles are uncontroversial as descriptions. The eliminativist must therefore reject not the principles but the description of the cases: she must say that in the duck-rabbit case what is \emph{given} is the same, and only the interpretation differs; that what is given to the person in the next room is the same as what is given to the person facing the cube, and only her access differs; and so on. The next section examines these redescriptions.

\section{The eliminations examined}

\subsection{Eliminating the sense}

The sense-datum theorist says that what is strictly given is a colored expanse, the same in the duck phase and the rabbit phase, and that \enquote{duck} and \enquote{rabbit} are interpretations added to it. Two objections to this redescription are standard and, in our view, decisive. First, the colored expanse is itself given as something: as a colored expanse, as having a shape, as being on the left. Removing the conceptual sense \enquote{duck} does not remove sense as such; it replaces a rich sense with a thin one. Second, as \citet{Sellars1956} argued, an item given under no sense could not stand in rational relations to beliefs, and so could not do the epistemic work for which the given was introduced. \citet{McDowell1994} draws the moral that the conceptual is already operative in experience. Our thesis does not require McDowell's conceptualism: the sense under which something is given may be non-conceptual, perspectival or practical. It requires only that there be one. The companion manuscript PHE-02-b develops the point that there is no pure datum.

Marion's proposal is more radical, because it does not locate the senseless given beneath interpretation but at the far end of the reduction. In \emph{Reduction and Givenness} he argues that Husserl reduced to objectness and Heidegger to being, and that a third reduction reaches givenness itself, answering to the pure form of the call \citep{Marion1989}. In \emph{Being Given} the saturated phenomenon gives more than any concept can contain \citep{Marion1997}. But excess of intuition over concept is not absence of sense. The saturated phenomenon is given as overwhelming, as unforeseeable, as exceeding: these are senses, and Marion's descriptions consist in specifying them. A phenomenon given as exceeding every determinate \enquote{as} is still given \emph{as} exceeding. Marion's analysis is therefore a contribution to the theory of the third place, showing that it can be filled by senses of excess, not an elimination of it.

\subsection{Eliminating the dative}

Heidegger's lecture \enquote{Time and Being} asks after the \enquote{it} in \enquote{it gives being} and \enquote{it gives time}, and answers with the event of appropriation, in which being is sent and withdraws \citep{Heidegger1969}. On one reading, this is a givenness without a dative. On a more careful reading, the event appropriates the human being to it, so that the dative is not removed but reconceived as belonging to the event rather than standing over against it. Either way, the giving at issue is not the givenness of a phenomenon, which is the relation phenomenology studies. If \enquote{it gives} means that there is something, the relation has become existence, which is one-place; and existence is not givenness, since much exists that is given to no one.

Marion's own theory keeps the dative. Its recipient, the \enquote{gifted} (l'adonné), receives itself from what it receives \citep{Marion1997}. That the recipient is constituted in the reception does not remove it from the relation; it gives it a particular kind of being. The dative place is required by the theory as well as by Principle~\ref{p:subject}.

\subsection{Eliminating the item}

\citet{Henry1963} distinguished the manifestation of the world, in which something appears at a distance in a horizon, from the manifestation of life, in which life feels itself without distance. In self-affection, there is no item other than the giving itself. One might take this to eliminate the first place, leaving a relation between a subject and a mode.

The reply is that the first place is not eliminated but filled by the subject. Self-affection is the case $G(s,s,F)$: life is given to life as feeling itself. A relation in which two places can be filled by the same value is not thereby a relation with fewer places, any more than self-love is a one-place relation. Henry's claim, which we need not dispute, is that this reflexive case is prior to the others. Priority is not elimination. Moreover, Henry's own descriptions differentiate the modes of self-affection, as suffering and as joy; these are values of the third place, and so the reflexive case is three-place as well.

\section{Objections}

\begin{objection}[Content]
The sense is just the content of the act, and contents are properties of the subject's state. So $G(x,s,F)$ reduces to the two-place relation between $x$ and the state $s$-having-$F$.
\end{objection}

\begin{reply}
Packing two places into one complex term does not reduce the arity of the relation; it hides it. The complex term \enquote{$s$ having $F$} has two independently variable components, and Principles~\ref{p:aspect} and~\ref{p:subject} vary them independently. Definition~\ref{d:elim} is formulated so that the reduction must be to a relation on two of the original places, not on two arbitrarily constructed terms.
\end{reply}

\begin{objection}[Time and circumstance]
Givenness has more than three places: times, conditions of illumination, the subject's location. If these are added, the claim of three places is false.
\end{objection}

\begin{reply}
The thesis is that givenness has at least the three places and that none of them can be eliminated. Times and circumstances are parameters on which each place may depend: the sense may be perspectival because the body is the zero-point of orientation (PHE-08), and it may change over time because things are given in adumbrations (PHE-03). These dependencies enrich the third place without competing with it.
\end{reply}

\begin{objection}[Externalism]
On disjunctivist and naïve realist views, perception is a two-place relation of acquaintance between a subject and an object, and the sense is at most a consequence of the relation, not a constituent of it \citep{Campbell2002}.
\end{objection}

\begin{reply}
Campbell's own view makes the relation three-place: the subject is related to the object from a standpoint, which fixes the way the object is presented. More generally, any view that admits Principle~\ref{p:aspect} must represent the difference between the duck phase and the rabbit phase somewhere; if not as a constituent of the relation, then as a variation in the relation that holds fixed both of its relata, which is what an irreducible third place is.
\end{reply}

\section{What the result does not establish}

The theorem concerns the logical form of givenness. It does not say what senses are, whether they are conceptual, or whether the subject is a transcendental ego or an embodied organism. Nor does it settle which givenness claims are true. It entails that every return to the things themselves returns to them as given to someone under some sense, which is why such a return is possible only in profiles and never all at once (PHE-03), and why perception, in which things are given in their most originary sense, can be the founding stratum of thought (PHE-14).

\section{Conclusion}

Givenness is a three-place relation. The same thing can be given to the same subject under different senses, the same thing under the same sense to one subject and not another, and different things to one subject under the same sense; so no place is fixed by the other two. The eliminations proposed by sense-datum theory, by Marion, by the late Heidegger and by Henry either reintroduce the place they remove, fill two places with the same value, or replace givenness with another relation.

\begin{thebibliography}{99}\small
\bibitem[Campbell(2002)]{Campbell2002} Campbell, J. (2002). \emph{Reference and Consciousness}. Oxford University Press.
\bibitem[Heidegger(1927)]{Heidegger1927} Heidegger, M. (1927). \emph{Sein und Zeit}. Niemeyer.
\bibitem[Heidegger(1969)]{Heidegger1969} Heidegger, M. (1969). Zeit und Sein. In \emph{Zur Sache des Denkens}. Niemeyer. (Lecture of 1962.)
\bibitem[Henry(1963)]{Henry1963} Henry, M. (1963). \emph{L'essence de la manifestation}. Presses Universitaires de France.
\bibitem[Husserl(1913)]{Husserl1913} Husserl, E. (1913). \emph{Ideen zu einer reinen Phänomenologie und phänomenologischen Philosophie. Erstes Buch}. Niemeyer.
\bibitem[Marion(1989)]{Marion1989} Marion, J.-L. (1989). \emph{Réduction et donation: recherches sur Husserl, Heidegger et la phénoménologie}. Presses Universitaires de France.
\bibitem[Marion(1997)]{Marion1997} Marion, J.-L. (1997). \emph{Étant donné: essai d'une phénoménologie de la donation}. Presses Universitaires de France.
\bibitem[McDowell(1994)]{McDowell1994} McDowell, J. (1994). \emph{Mind and World}. Harvard University Press.
\bibitem[Sellars(1956)]{Sellars1956} Sellars, W. (1956). Empiricism and the philosophy of mind. In H. Feigl \& M. Scriven (Eds.), \emph{Minnesota Studies in the Philosophy of Science}, Vol.~1 (pp.~253--329). University of Minnesota Press.
\bibitem[Wittgenstein(1953)]{Wittgenstein1953} Wittgenstein, L. (1953). \emph{Philosophical Investigations}. Blackwell.
\end{thebibliography}

\end{document}
